\documentclass[11pt]{article}
\usepackage[a4paper,margin=1in]{geometry}
\usepackage{amsmath,amssymb,amsthm,mathtools}
\usepackage{graphicx}
\usepackage{microtype}
\usepackage[hidelinks]{hyperref}

\newtheorem{theorem}{Theorem}
\newtheorem{lemma}[theorem]{Lemma}
\newtheorem{corollary}[theorem]{Corollary}
\newcommand{\C}{\mathbb C}
\newcommand{\id}{\operatorname{id}}
\newcommand{\rank}{\operatorname{rank}}

\title{The Near-Critical Sample Tradeoff for Uniform-Norm Projections\\
\large Endpoint resolution and quantitative bounds for arXiv:2401.02220}
\author{Automated Mathematics Run \texttt{fa\_banach\_001}, lane 0}
\date{11 August 2026}

\begin{document}
\maketitle

\begin{center}
\fbox{\parbox{0.91\linewidth}{\textbf{Status.}
Candidate partial result, likely valid.  The case of one extra sample
(equivalently $\alpha=0$) is completely resolved up to absolute constants.
For $0<\alpha<1$ the packet gives new explicit brackets, but the full
asymptotic tradeoff remains open.}}
\end{center}

\section{Source problem}

Krieg, Pozharska, M.~Ullrich, and T.~Ullrich prove that every
$n$-dimensional subspace of $B(D)$ admits a sampling projection using $2n$
function values and having norm $O(\sqrt n)$ \cite{KPUU}.  At the other
endpoint, $n$ samples require norm of order $n$ in the worst case.  They ask
for the smallest universal norm when only $n+n^\alpha$ values are allowed,
$0\leq\alpha<1$.  The exact source passage is reproduced below.

\begin{figure}[ht]
  \centering
  \includegraphics[width=0.96\linewidth]{figures/open_problem_crop.png}
  \caption{The open problem on page 2 of arXiv:2401.02220v2.}
\end{figure}

To make ``smallest universal norm'' precise, let
\[
 \Pi(n,m)=\sup_{D,V}\ \inf_P \|P\|_{B(D)\to B(D)},
\]
where the supremum is over all sets $D$ and all complex $n$-dimensional
subspaces $V\subset B(D)$, and the infimum is over sampling projections
\[
 Pf=\sum_{i=1}^r f(x_i)\varphi_i,
 \qquad r\leq m,\quad x_i\in D,\quad \varphi_i\in V,
\]
onto $V$.  Repeated sample points may be merged, so they never improve the
infimum.

\section{Result}

\begin{theorem}[Near-critical sampling bracket]\label{thm:main}
For $n\geq2$ and $1\leq k\leq n$,
\begin{equation}\label{eq:lower}
 \Pi(n,n+k)\geq
 \max\left\{\frac{n}{\sqrt{n+k}},
              \left\lfloor\frac{n}{k+1}\right\rfloor\right\}.
\end{equation}
For every such $k$ one has $\Pi(n,n+k)\leq n+1$.  If $k\geq2$, then also
\begin{align}
 \Pi(n,n+k)
 &\leq \sqrt{n+2}\,
 \frac{\sqrt{n+k}+\sqrt{n+1}}
      {\sqrt{n+k}-\sqrt{n+1}}                                      \label{eq:upper-exact}\\
 &\leq \frac{4(n+k)\sqrt{n+2}}{k-1}.                               \label{eq:upper-simple}
\end{align}
Consequently the better of $n+1$ and \eqref{eq:upper-exact} may always be
used.
\end{theorem}

\begin{corollary}[Power oversampling]\label{cor:alpha}
Let $k_n=\lceil n^\alpha\rceil$.  For $\alpha=0$,
\[
 \Pi(n,n+1)=\Theta(n).
\]
Thus the first previously open point of the tradeoff curve is fully
resolved.  For fixed $0<\alpha<1$,
\begin{equation}\label{eq:power-bracket}
 n^{\max\{1/2,,1-\alpha\}}
 \ \lesssim\ \Pi(n,n+k_n)
 \ \lesssim\
 n^{\min\{1,,3/2-\alpha\}}.
\end{equation}
In particular, the upper bound becomes nontrivial already for
$\alpha>1/2$.  The unweighted estimate stated in Remark 7 of the source
paper becomes nontrivial only for $\alpha>2/3$.
\end{corollary}

\section{Proof intuition}

The two lower bounds detect different obstructions.  Linear functions on
the Euclidean sphere force every $m$-sample reconstruction of the identity
to have norm at least $n/\sqrt m$; this follows just by taking a trace and
then averaging over random phases.  To see the instability at the critical
sample count, use the zero-sum hyperplane
\[
 H_q=\{z\in\C^{q+1}:z_0+\cdots+z_q=0\}.
\]
Sampling exactly $q$ of its $q+1$ coordinates forces the omitted coordinate
to be minus the sum of the sampled values, and hence forces projection norm
$q$.  Taking $k+1$ disjoint blocks guarantees that one block receives no
extra sample, which gives the second term in \eqref{eq:lower}.

The upper bound uses positivity in a different way.  Adjoin the constant
function to $V$.  A Kiefer--Wolfowitz design gives a probability measure for
which the uniform norm is bounded by $\sqrt{n+1}$ times the $L^2$ norm.
Weighted BSS sparsification reduces this measure to $n+k$ atoms while
losing a factor of order $n/k$ in the $L^2$ comparison.  The adjoined
constant controls the total sparse weight.  Weighted least squares is then
a sampling projection with the same order of norm.  This last constant-
function device is essential: without it, the sparse weights need not have
controlled total mass.

\section{Lower bounds}

\begin{lemma}[Euclidean trace obstruction]\label{lem:euclidean}
For every $m\geq n$, $\Pi(n,m)\geq n/\sqrt m$.
\end{lemma}

\begin{proof}
Let $D$ be the Euclidean unit sphere of $\C^n$, and identify $\C^n$ with the
space $V$ of linear functions $f_z(x)=\langle z,x\rangle$ on $D$.  This is
an isometry from $\ell_2^n$ into $B(D)$.

Suppose a sampling projection uses distinct points $x_1,\dots,x_r$, where
$r\leq m$.  Write its output functions as $f_{a_1},\dots,f_{a_r}$.  Since
the operator is the identity on $V$,
\begin{equation}\label{eq:identity-decomp}
 z=\sum_{i=1}^r\langle z,x_i\rangle a_i
 \qquad(z\in\C^n).
\end{equation}
Taking traces and then absolute values gives
\[
 n\leq\sum_{i=1}^r\|a_i\|_2
 \leq\sqrt r\left(\sum_{i=1}^r\|a_i\|_2^2\right)^{1/2}.
\]
For independent uniform phases $\varepsilon_i$ on the unit circle,
\[
 \mathbb E\left\|\sum_i\varepsilon_i a_i\right\|_2^2
 =\sum_i\|a_i\|_2^2.
\]
Because an arbitrary bounded function on $D$ may take the prescribed values
$\varepsilon_i$ at the distinct sample points, the square root on the
right is at most the projection norm.  Hence
$n\leq\sqrt r\|P\|\leq\sqrt m\|P\|$.
\end{proof}

\begin{lemma}[Block-hyperplane obstruction]\label{lem:block}
For $1\leq k\leq n$,
\[
 \Pi(n,n+k)\geq\left\lfloor\frac{n}{k+1}\right\rfloor.
\]
\end{lemma}

\begin{proof}
Put $b=k+1$ and $q=\lfloor n/b\rfloor$.  On a set of $q+1$ points let
\[
 H_q=\{z\in\C^{q+1}:\textstyle\sum_{i=0}^qz_i=0\},
\]
with its inherited uniform norm.  Form the $\ell_\infty$-direct sum of $b$
disjoint copies of $H_q$ and add $n-bq$ one-dimensional coordinate blocks.
The resulting function space $V\subset B(D)$ has dimension $n$.

Any sampling projection onto $V$ must use at least $q$ distinct points in
each $H_q$ block and at least one in each one-dimensional block: otherwise
the restriction of the evaluation map to that summand is not injective.
These compulsory samples account for $n$ points.  There are only $k=b-1$
extra points, so at least one $H_q$ block is sampled at exactly $q$ of its
$q+1$ coordinates.

Restrict inputs to that block and take the corresponding output component.
This gives a $q$-sample projection from $\ell_\infty^{q+1}$ onto $H_q$.
The $q$ sampled-coordinate functionals restrict to an isomorphism on
$H_q$, so its left inverse is unique: it copies the sampled values and puts
minus their sum in the omitted coordinate.  Its norm is exactly $q$.
Therefore the original projection also has norm at least $q$.
\end{proof}

Lemmas \ref{lem:euclidean} and \ref{lem:block} prove
\eqref{eq:lower}.  Notice that the second construction is what survives
one extra sample: a single zero-sum block becomes harmless after sampling
all its coordinates, but $k+1$ blocks ensure that $k$ extra samples cannot
repair all of them.

\section{Upper bound}

We use the following finite-frame sparsification consequence of
Bartel--Sch\"afer--Ullrich \cite{BSU}.  It is their preconditioned complex
BSS estimate, specialized to a finite-dimensional subspace.

\begin{lemma}[Weighted near-minimal discretization]\label{lem:bss}
Let $W$ be a $d$-dimensional subspace of $L^2(\rho)$, where $\rho$ is a
finitely supported probability measure.  For every $c>1$ there are at most
$\lceil cd\rceil$ atoms $x_j$ and weights $\lambda_j\geq0$ such that
\begin{equation}\label{eq:bss}
 \|w\|_{L^2(\rho)}^2
 \leq\sum_j\lambda_j|w(x_j)|^2
 \leq\gamma_c\|w\|_{L^2(\rho)}^2,
 \qquad
 \gamma_c=\left(\frac{\sqrt c+1}{\sqrt c-1}\right)^2.
\end{equation}
\end{lemma}

\begin{proof}[Proof of the upper bounds in Theorem \ref{thm:main}]
The $n+1$ bound is the sampling version of Auerbach's lemma recalled in
Remark 4 of \cite{KPUU}; extra samples may simply be ignored.

Assume $m=n+k\geq n+2$.  Set $W=\operatorname{span}(V\cup\{1\})$ and
$d=\dim W\leq n+1$.  Proposition 8 of \cite{KPUU}, applied with its
parameter $\varepsilon=1$, gives a finitely supported probability measure
$\rho$ such that
\begin{equation}\label{eq:design}
 \|w\|_\infty\leq\sqrt{d+1}\,\|w\|_{L^2(\rho)}
 \qquad(w\in W).
\end{equation}
Apply Lemma \ref{lem:bss} with $c=m/d>1$.  Let
\(\|f\|_\nu^2=\sum_j\lambda_j|f(x_j)|^2\).
Since $1\in W$, the upper half of \eqref{eq:bss} gives
\begin{equation}\label{eq:mass}
 \nu(D)=\sum_j\lambda_j=\|1\|_\nu^2\leq\gamma_c.
\end{equation}

Let $A_\nu f$ be the weighted least-squares minimizer over $V$.  The lower
half of \eqref{eq:bss} makes this minimizer unique; it is linear, depends on
at most $m$ values of $f$, and is the identity on $V$.  Thus it is a
sampling projection.  Orthogonal projection in the $\nu$ inner product,
\eqref{eq:design}, \eqref{eq:bss}, and \eqref{eq:mass} give
\begin{align*}
 \|A_\nu f\|_\infty
 &\leq\sqrt{d+1}\|A_\nu f\|_{L^2(\rho)}
 \leq\sqrt{d+1}\|A_\nu f\|_\nu\\
 &\leq\sqrt{d+1}\|f\|_\nu
 \leq\sqrt{(d+1)\gamma_c}\|f\|_\infty.
\end{align*}
Therefore
\[
 \|A_\nu\|
 \leq\sqrt{d+1}\frac{\sqrt m+\sqrt d}{\sqrt m-\sqrt d}
 \leq\sqrt{n+2}\frac{\sqrt m+\sqrt{n+1}}
                         {\sqrt m-\sqrt{n+1}},
\]
which is \eqref{eq:upper-exact}.  Finally,
\[
 \frac{\sqrt m+\sqrt d}{\sqrt m-\sqrt d}
 =\frac{(\sqrt m+\sqrt d)^2}{m-d}
 \leq\frac{4m}{m-d}
 \leq\frac{4(n+k)}{k-1},
\]
proving \eqref{eq:upper-simple}.
\end{proof}

Corollary \ref{cor:alpha} follows immediately by putting
$k=\lceil n^\alpha\rceil$ in Theorem \ref{thm:main}.

\section{Upgrade attempts, verification, and remaining gap}

The proof was pushed beyond the first partial observation in four ways.
First, the unweighted sparsification used in the source was replaced by the
weighted preconditioned BSS estimate, and the constant function was adjoined
to control total mass.  Second, two independent lower witnesses were made
explicit rather than importing the paper's sharpness assertion.  Third,
tensor products of zero-sum hyperplanes were tested as possible stronger
obstructions.  Exact small linear programs give best norms $3,2$ for
$H_2\otimes H_2$ with $4,5$ samples and $4,3,2$ for
$H_2\otimes H_3$ with $6,7,8$ samples; the rapid improvement rules out the
naive tensor route.  Fourth, an interpolation-plus-rank-one-correction
decomposition was examined.  Although each correction can be chosen to
annihilate $V$ and hence preserve the projection property, no selection
principle controls the uniform operator norm of a partial correction sum;
the necessary cancellations are global in the output rows.

The script \texttt{code/verify\_finite\_witnesses.py} checks the unique
$q$-sample projection onto $H_q$ for $1\leq q\leq12$, verifies its norm is
$q$, and checks the block-count formulas and scalar bounds in the theorem.
These computations are consistency checks and are not used in the proof.

For $0<\alpha<1$, the exponent gap in \eqref{eq:power-bracket} remains.
Closing it appears to require either a sampling projection construction that
uses cancellations unavailable to positive weighted least squares, or a
single family of subspaces whose instability persists quantitatively under
many extra coordinates.  The block construction only charges one extra
sample to repair one hard block, and weighted spectral sparsification
intrinsically pays the $(c-1)^{-1}$ loss.

\section{Novelty assessment}

The exact title, arXiv identifier, open-problem phrase, ``$n+k$ sampling
projection'', ``near-critical oversampling'', and related frame-
sparsification terms were searched through 11 August 2026.  The October 2025
version of the source still poses the problem.  A March 2026 survey by Dai,
Kosov, and Temlyakov \cite{DKT} gives a weighted norming-set estimate of
order $\sqrt n/(c-1)$ by adjoining the constant function, but does not state
the sampling-projection tradeoff, the block lower bound, or the one-extra-
sample resolution.  Thus the upper argument should be viewed as a short
projection consequence of known ingredients, not as a wholly new
sparsification theorem.  No source located in the bounded search contains
the block-hyperplane lower bound or the conclusion
$\Pi(n,n+1)=\Theta(n)$.  Novelty confidence for those additions is moderate,
not definitive.

\paragraph{Human-review recommendation.}
Review as a likely-valid partial result with a fully resolved endpoint.  The
highest-value checks are the uniqueness argument on an exactly sampled
zero-sum block, the count forcing one such block, and the use of the constant
function in \eqref{eq:mass}.  The packet makes no claim to resolve the
interior range $0<\alpha<1$.

\begin{thebibliography}{9}
\bibitem{KPUU}
D.~Krieg, K.~Pozharska, M.~Ullrich, and T.~Ullrich,
\emph{Sampling projections in the uniform norm},
J. Math. Anal. Appl. \textbf{553} (2026), Paper 129873;
arXiv:2401.02220.

\bibitem{BSU}
F.~Bartel, M.~Sch\"afer, and T.~Ullrich,
\emph{Constructive subsampling of finite frames with applications in
optimal function recovery},
Appl. Comput. Harmon. Anal. \textbf{65} (2023), 209--248;
arXiv:2202.12625.

\bibitem{DKT}
F.~Dai, E.~Kosov, and V.~N.~Temlyakov,
\emph{A survey of sampling discretization of integral and uniform norms},
arXiv:2603.02459 (2026).
\end{thebibliography}

\end{document}
